% !TeX root = ../概率论与数理统计笔记.tex
% 由主文档通过 \input 载入；请编译主文档。
\setcounter{chapter}{11} % 保留原书章号。
\chapter{点估计及其评价}
\label{chap:ps-point-estimation}

本章将未知参数估计为一个由样本计算的数值，并建立比较估计方法的标准。
先修内容为抽样分布与似然；涉及最优无偏估计时使用第\ref{chap:ps-information}章，涉及大样本性质时使用第\ref{chap:ps-limits}章。

\section{估计目标、损失与风险}
\label{sec:ps-estimation-criteria}
给定参数函数 $g(\theta)$，以 $\widehat g=\widehat g(\mathbf X)$ 估计它；定义偏差、方差与\emph{均方误差}（mean squared error, MSE）。
在二阶矩存在时推导均方误差等于方差加偏差平方，说明无偏性并不自动意味着误差最小。
引入\emph{损失函数}（loss function）与\emph{风险函数}（risk function），明确比较范围及其对评价结果的影响。

\section{矩估计}
\label{sec:ps-method-moments}
介绍\emph{矩估计}（method of moments）：用样本矩替代总体矩并求解参数方程。
说明所用矩的存在、参数可识别性、解的存在唯一性与参数空间约束。
通过单参数和双参数模型比较计算简便性与统计效率，不把方程有解等同于解一定合理。

\section{最大似然估计}
\label{sec:ps-maximum-likelihood}
定义\emph{最大似然估计}（maximum likelihood estimation, MLE）为在 $\Theta$ 上最大化给定观测的似然。
讨论对数似然、边界解、多解、上确界不取到以及参数相关支撑的情形。
得分方程只用于寻找适当内点候选，最终需检查可行域与整体最大性；不变性在正确处理参数映射与解集后表述。

\section{无偏性与最小方差无偏估计}
\label{sec:ps-unbiased-estimation}
通过样本均值、样本方差与均匀分布端点估计比较无偏修正。
运用充分统计量、条件化与完备性构造\emph{一致最小方差无偏估计}（uniformly minimum-variance unbiased estimator, UMVUE）。
明确“最优”限于对所有参数均无偏的估计量类，不与所有可能损失下的最优性混同。

\section{相合性与渐近性质}
\label{sec:ps-estimation-asymptotics}
定义\emph{相合性}（consistency，又称一致性）与渐近正态性，区分它们与有限样本无偏性。
用大数定律与连续映射证明典型矩估计的相合性。
最大似然估计的渐近结论列出可识别性、内点真参数、适当光滑性和信息非退化等条件，并通过边界或不规则模型说明结论并非无条件成立。

\section{例题、习题与学习检查}
规划伯努利、泊松、正态和均匀端点模型的矩估计与最大似然估计；比较均匀端点估计的偏差和均方误差。
习题既包含计算，也包含“无偏但不相合”“有偏但相合”的例子与辨析。
学习检查要求分别报告估计目标、估计量公式、观测后的估计值和评价标准。






